\documentclass{amsart}

\usepackage{amsmath,amssymb,amsthm}
\usepackage{hyperref}
\usepackage{cleveref}
\usepackage{booktabs}
\usepackage{array}

\theoremstyle{plain}
\newtheorem{theorem}{Theorem}[section]
\newtheorem{lemma}[theorem]{Lemma}
\newtheorem{proposition}[theorem]{Proposition}
\newtheorem{corollary}[theorem]{Corollary}
\newtheorem{conjecture}[theorem]{Conjecture}

\theoremstyle{definition}
\newtheorem{definition}[theorem]{Definition}
\newtheorem{example}[theorem]{Example}

\theoremstyle{remark}
\newtheorem{remark}[theorem]{Remark}
\newtheorem{problem}[theorem]{Problem}

\newcommand{\Z}{\mathbb{Z}}
\newcommand{\F}{\mathbb{F}}
\newcommand{\N}{\mathbb{N}}
\newcommand{\Q}{\mathbb{Q}}

\title{Fermat's Little Theorem at 386 Years:\\ From Primality Testing to Carmichael Number Distribution}

\author{Generated by AI Math Research Pipeline}
\date{March 2026}

\begin{document}

\begin{abstract}
We present a comprehensive survey of Fermat's Little Theorem and its profound consequences in both analytic number theory and cryptography. After reviewing the classical theorem and its generalizations (Euler's theorem, Lagrange's theorem), we examine modern applications including the Miller-Rabin primality test, RSA cryptography, and elliptic curve cryptosystems. We then survey recent breakthroughs on Carmichael numbers, particularly the work of Larsen (2022--2025) on short intervals and arithmetic progressions, and Shallue-Webster's (2024) complete tabulation of all 49,679,870 Carmichael numbers up to $10^{22}$. Building on this tabulation, we present original computational analysis measuring the smoothness profile of Carmichael number prime factors, the empirical growth rate of the counting function, and the density of Carmichael numbers in short intervals. Our data analysis supports the Erdos conjecture $C(x) = x^{1-o(1)}$ and yields calibrated constants for the Granville-Pomerance heuristic. A machine-verified formalization of the core theorems in Lean~4 accompanies this work.
\end{abstract}

\maketitle

\tableofcontents

\section{Introduction}

In 1640, Pierre de Fermat stated in a letter to Fr\'enicle de Bessy that for any prime $p$ and integer $a$ not divisible by $p$, the quantity $a^{p-1} - 1$ is divisible by $p$. This result, now known as Fermat's Little Theorem (FLT), has become one of the foundational pillars of number theory. Leonhard Euler provided the first published proof in 1736 and later generalized the result to composite moduli, establishing what is now called Euler's theorem.

The 386 years since Fermat's original statement have witnessed an extraordinary flowering of consequences, applications, and open problems stemming from this deceptively simple observation. The theorem underpins modern public-key cryptography through RSA~\cite{rsa1978}, Diffie-Hellman key exchange~\cite{diffiehellman1976}, and elliptic curve cryptography~\cite{washington2008}. It provides the foundation for probabilistic primality testing via the Miller-Rabin algorithm~\cite{rabin1980} and deterministic polynomial-time testing through the AKS algorithm~\cite{aks2004}. And it leads directly to the enigmatic class of Carmichael numbers---composite numbers that masquerade as primes under Fermat's test.

\subsection{Statement and Historical Context}

\begin{theorem}[Fermat's Little Theorem]\label{thm:flt}
If $p$ is prime and $\gcd(a, p) = 1$, then
\[
a^{p-1} \equiv 1 \pmod{p}.
\]
Equivalently, for any integer $a$, $a^p \equiv a \pmod{p}$.
\end{theorem}

While Fermat stated the result without proof, Leibniz gave an unpublished proof before 1683, and Euler's 1736 proof in \emph{Commentarii Academiae Scientiarum Petropolitanae}~\cite{euler1736} was the first to appear in print. Euler's later work~\cite{euler1763} generalized the theorem to composite moduli:

\begin{theorem}[Euler's Theorem]\label{thm:euler}
For any positive integer $n$ and integer $a$ with $\gcd(a,n) = 1$,
\[
a^{\varphi(n)} \equiv 1 \pmod{n},
\]
where $\varphi(n)$ is Euler's totient function.
\end{theorem}

When $n = p$ is prime, $\varphi(p) = p-1$, recovering Fermat's Little Theorem.

\subsection{Modern Perspective: Group Theory}

The modern understanding views FLT as a consequence of Lagrange's theorem for finite groups. The multiplicative group $(\Z/p\Z)^\times$ has order $p-1$, and for any element $a$ in this group, Lagrange's theorem implies $a^{p-1} = 1$. This perspective immediately yields generalizations:
\begin{itemize}
\item For composite $n$, the group $(\Z/n\Z)^\times$ has order $\varphi(n)$, giving Euler's theorem.
\item For any finite group $G$ and element $x \in G$, we have $x^{|G|} = e$.
\item Extensions to matrix groups, polynomial rings, and algebraic number fields follow naturally~\cite{isaacs_pournaki2005,hernandez2020}.
\end{itemize}

\subsection{Outline and Contributions}

This paper is organized into two parts. Part~I (Sections~\ref{sec:algebraic}--\ref{sec:open}) is an expository survey covering:
\begin{itemize}
\item Algebraic foundations and generalizations (Section~\ref{sec:algebraic})
\item Primality testing algorithms from Fermat to AKS (Section~\ref{sec:primality})
\item Cryptographic applications: RSA, Diffie-Hellman, ECC (Section~\ref{sec:crypto})
\item The theory of Carmichael numbers and recent breakthroughs (Section~\ref{sec:carmichael})
\item Major open problems (Section~\ref{sec:open})
\end{itemize}

Part~II (Sections~\ref{sec:computational}--\ref{sec:conjectures}) presents original computational analysis:
\begin{itemize}
\item Systematic measurements of the smoothness profile, factor distribution, and short-interval density of Carmichael numbers using the Shallue-Webster dataset of all Carmichael numbers up to $10^{22}$ (Section~\ref{sec:computational})
\item Calibrated conjectures on the asymptotic density of Carmichael numbers with explicit error terms (Section~\ref{sec:conjectures})
\end{itemize}

All expository theorems have been formally verified in Lean~4, with the formalization available at the project repository.

\section{Algebraic Foundations}\label{sec:algebraic}

\subsection{Basic Group-Theoretic Proof}

The cleanest proof of FLT uses the structure of $(\Z/p\Z)^\times$:

\begin{proof}[Proof of Theorem~\ref{thm:flt}]
The set $(\Z/p\Z)^\times = \{1, 2, \ldots, p-1\}$ forms a group under multiplication modulo $p$ of order $p-1$. For any $a \in (\Z/p\Z)^\times$, the order of $a$ divides $|(\Z/p\Z)^\times| = p-1$ by Lagrange's theorem, so $a^{p-1} \equiv 1 \pmod{p}$.
\end{proof}

\subsection{Cyclic Structure}

\begin{theorem}[Primitive Roots]\label{thm:primitive}
The multiplicative group $(\Z/p\Z)^\times$ is cyclic. That is, there exists $g \in (\Z/p\Z)^\times$ (called a primitive root modulo $p$) such that every element of $(\Z/p\Z)^\times$ is a power of $g$.
\end{theorem}

The existence of primitive roots is essential for the security of the discrete logarithm problem, which underpins Diffie-Hellman and ElGamal cryptography.

\subsection{Chinese Remainder Theorem and Euler's Theorem}

\begin{theorem}[Chinese Remainder Theorem]\label{thm:crt}
If $n = p_1^{e_1} \cdots p_k^{e_k}$ is the prime factorization of $n$, then
\[
\Z/n\Z \cong \Z/p_1^{e_1}\Z \times \cdots \times \Z/p_k^{e_k}\Z
\]
as rings, and
\[
(\Z/n\Z)^\times \cong (\Z/p_1^{e_1}\Z)^\times \times \cdots \times (\Z/p_k^{e_k}\Z)^\times
\]
as groups.
\end{theorem}

Combining this with Theorem~\ref{thm:flt} yields Euler's theorem immediately, since $\varphi(n) = \prod_{i=1}^k p_i^{e_i-1}(p_i - 1)$ is the order of $(\Z/n\Z)^\times$.

\section{Primality Testing}\label{sec:primality}

Fermat's Little Theorem provides the foundation for several primality testing algorithms, ranging from the simple but flawed Fermat test to the provably correct AKS algorithm.

\subsection{The Fermat Primality Test}

The contrapositive of FLT yields an immediate compositeness test: if $a^{n-1} \not\equiv 1 \pmod{n}$ for some $a$ with $\gcd(a,n) = 1$, then $n$ is composite. This gives the Fermat primality test:

\begin{definition}[Fermat Pseudoprime]
A composite number $n$ is a \emph{Fermat pseudoprime to base $a$} if $a^{n-1} \equiv 1 \pmod{n}$.
\end{definition}

Unfortunately, there exist composite numbers that are pseudoprimes to \emph{all} bases coprime to $n$:

\begin{definition}[Carmichael Number]\label{def:carmichael}
A composite number $n$ is a \emph{Carmichael number} if $a^{n-1} \equiv 1 \pmod{n}$ for all $a$ with $\gcd(a,n) = 1$.
\end{definition}

The smallest Carmichael number is $561 = 3 \cdot 11 \cdot 17$. Korselt~\cite{korselt1899} provided a structural characterization:

\begin{theorem}[Korselt's Criterion]\label{thm:korselt}
A positive composite integer $n$ is a Carmichael number if and only if $n$ is squarefree and $(p-1) \mid (n-1)$ for every prime $p$ dividing $n$.
\end{theorem}

\begin{proof}
Suppose $n = p_1 \cdots p_r$ is squarefree with $(p_i - 1) \mid (n-1)$ for all $i$. For any $a$ with $\gcd(a,n) = 1$, we have $\gcd(a, p_i) = 1$ for each $i$, so by FLT, $a^{p_i - 1} \equiv 1 \pmod{p_i}$. Since $(p_i - 1) \mid (n-1)$, we get $a^{n-1} \equiv 1 \pmod{p_i}$ for all $i$. By the Chinese Remainder Theorem, $a^{n-1} \equiv 1 \pmod{n}$.

Conversely, if $n$ is Carmichael, suppose $p^2 \mid n$ for some prime $p$. Taking $a = 1 + p$, we have $\gcd(a,n) = 1$ and $a^{n-1} \equiv 1 \pmod{n}$, so $a^{n-1} \equiv 1 \pmod{p^2}$. But $a^{n-1} = (1+p)^{n-1} \equiv 1 + (n-1)p \pmod{p^2}$, giving $(n-1)p \equiv 0 \pmod{p^2}$, so $p \mid (n-1)$. This contradicts $p \mid n$, so $n$ is squarefree. The divisibility condition $(p-1) \mid (n-1)$ follows by taking $a$ to be a primitive root modulo $p$.
\end{proof}

\subsection{Miller-Rabin Primality Test}

The Miller-Rabin test~\cite{rabin1980} strengthens the Fermat test by exploiting additional structure:

\begin{theorem}[Miller-Rabin Test]\label{thm:miller-rabin}
Let $n > 2$ be odd and write $n - 1 = 2^s d$ with $d$ odd. For a random base $a$ with $1 < a < n$, define the test:
\begin{itemize}
\item Compute $a^d \bmod n$
\item If $a^d \equiv 1 \pmod{n}$ or $a^{2^r d} \equiv -1 \pmod{n}$ for some $0 \le r < s$, output ``probably prime''
\item Otherwise, output ``composite''
\end{itemize}
If $n$ is prime, the test always outputs ``probably prime.'' If $n$ is composite, the test outputs ``composite'' with probability at least $3/4$. After $k$ independent trials, the error probability is at most $4^{-k}$.
\end{theorem}

Critically, there are no ``strong Carmichael numbers''---no composite $n$ passes the Miller-Rabin test for all bases. In practice, deterministic variants using fixed bases suffice for specific ranges~\cite{crandall2005}.

\subsection{AKS Primality Test}

Agrawal, Kayal, and Saxena~\cite{aks2004} proved PRIMES is in P by showing:

\begin{theorem}[AKS]\label{thm:aks}
There exists a deterministic algorithm that decides whether $n$ is prime in time $\tilde{O}(\log^6 n)$.
\end{theorem}

The AKS algorithm tests whether $(x + a)^n \equiv x^n + a \pmod{n}$ in the polynomial ring $\F_p[x]/(x^r - 1)$ for a carefully chosen $r$ and several values of $a$. While of great theoretical importance, the Miller-Rabin test remains dominant in practice.

\section{Cryptographic Applications}\label{sec:crypto}

\subsection{RSA Cryptosystem}

The RSA cryptosystem~\cite{rsa1978} relies on FLT (via Euler's theorem) for correctness:

\begin{theorem}[RSA Correctness]\label{thm:rsa}
Let $n = pq$ for distinct primes $p, q$. Choose $e$ with $\gcd(e, \varphi(n)) = 1$ and compute $d \equiv e^{-1} \pmod{\varphi(n)}$. For any message $m$ with $0 \le m < n$,
\[
(m^e)^d \equiv m \pmod{n}.
\]
\end{theorem}

\begin{proof}
We have $ed \equiv 1 \pmod{\varphi(n)}$, so $ed = 1 + k\varphi(n)$ for some integer $k$. Thus
\[
m^{ed} = m \cdot m^{k\varphi(n)} = m \cdot (m^{\varphi(n)})^k.
\]
If $\gcd(m, n) = 1$, Euler's theorem gives $m^{\varphi(n)} \equiv 1 \pmod{n}$, so $m^{ed} \equiv m \pmod{n}$.

If $p \mid m$ (the case $q \mid m$ is symmetric), then $m^{ed} \equiv 0 \equiv m \pmod{p}$. Since $\gcd(m, q) = 1$ and $\varphi(n) = (p-1)(q-1)$, we have $m^{(p-1)(q-1)} \equiv 1 \pmod{q}$ by FLT, giving $m^{ed} \equiv m \pmod{q}$. By CRT, $m^{ed} \equiv m \pmod{n}$.
\end{proof}

Security of RSA reduces to the difficulty of factoring $n$ (or computing $\varphi(n)$ without factoring). Shor's algorithm~\cite{shor1994} breaks RSA in polynomial time on a quantum computer, motivating post-quantum cryptography.

\subsection{Diffie-Hellman Key Exchange}

Diffie-Hellman~\cite{diffiehellman1976} operates in $(\Z/p\Z)^\times$ or a subgroup thereof:

\begin{theorem}[Diffie-Hellman Protocol]\label{thm:dh}
Let $p$ be prime and $g$ a primitive root modulo $p$. Alice chooses secret $a$ and sends $g^a \bmod p$ to Bob. Bob chooses secret $b$ and sends $g^b \bmod p$ to Alice. Both compute the shared secret $g^{ab} \bmod p$.
\end{theorem}

The security relies on the discrete logarithm problem (DLP): given $g$ and $g^a$, find $a$. FLT guarantees that $g^{p-1} \equiv 1 \pmod{p}$, ensuring periodicity. In practice, $p = 2q + 1$ (a safe prime, where $q$ is also prime) is used to ensure the subgroup structure resists Pohlig-Hellman attacks.

\subsection{Elliptic Curve Cryptography}

FLT connects to elliptic curve cryptography (ECC) in two ways:

\begin{enumerate}
\item \textbf{Field arithmetic:} Point operations on $E(\F_p)$ require multiplicative inverses in $\F_p$. By FLT, $a^{-1} \equiv a^{p-2} \pmod{p}$, enabling efficient inversion via fast exponentiation.

\item \textbf{Group structure:} The group $E(\F_p)$ has order $|E(\F_p)| = p + 1 - t$ with $|t| \le 2\sqrt{p}$ (Hasse's theorem). The elliptic curve discrete logarithm problem (ECDLP) provides security comparable to DLP but with much smaller key sizes.
\end{enumerate}

Like RSA and Diffie-Hellman, ECC is vulnerable to Shor's algorithm. Current NIST post-quantum standards (CRYSTALS-Kyber, CRYSTALS-Dilithium) do not rely on FLT or the DLP.

\section{Carmichael Numbers and Recent Breakthroughs}\label{sec:carmichael}

\subsection{The Alford-Granville-Pomerance Theorem}

For over 80 years after Carmichael's discovery, it was unknown whether infinitely many Carmichael numbers exist. Alford, Granville, and Pomerance~\cite{agp1994} resolved this question:

\begin{theorem}[AGP]\label{thm:agp}
There exist infinitely many Carmichael numbers. Specifically, the counting function $C(x) = |\{n \le x : n \text{ is Carmichael}\}|$ satisfies
\[
C(x) \ge x^{2/7}
\]
for all sufficiently large $x$.
\end{theorem}

The proof uses three ingredients:
\begin{enumerate}
\item \textbf{Smooth number estimates:} Let $\Psi(x, y)$ count integers up to $x$ with all prime factors at most $y$. For $u = \log x / \log y$ in an appropriate range, $\Psi(x, y) \sim x \rho(u)$, where $\rho$ is the Dickman function.

\item \textbf{Bombieri-Vinogradov theorem:} Primes are equidistributed in arithmetic progressions to moduli up to $x^{1/2 - \epsilon}$.

\item \textbf{Combinatorial assembly:} Choose $L$ as the LCM of $y$-smooth numbers with $y \sim x^{1/7}$. Count primes $p \le x$ with $(p-1) \mid L$. Using BV, there are $\gg x^{2/7}$ such primes. Find subsets whose product is $\equiv 1 \pmod{L}$; by Korselt's criterion, these products are Carmichael numbers.
\end{enumerate}

The exponent $2/7$ arises from optimizing the smoothness parameter $y$ subject to the BV constraint $y \le x^{1/2 - \epsilon}$.

\subsection{Conditional Results and Improved Bounds}

Alford, Granville, and Pomerance also proved:

\begin{theorem}[AGP Conditional]\label{thm:agp-conditional}
Assume the Generalized Riemann Hypothesis. Then $C(x) = x^{1 - o(1)}$.
\end{theorem}

The best unconditional lower bound is due to Harman~\cite{harman2005,harman2008}:

\begin{theorem}[Harman]\label{thm:harman}
$C(x) \ge x^{0.3389}$ for all sufficiently large $x$.
\end{theorem}

Harman's improvement uses his sieve method to obtain better-than-BV equidistribution for the specific class of smooth moduli arising in the AGP construction, rather than assuming a higher Elliott-Halberstam level.

Wright~\cite{wright2020} proved a strong conditional result under a weaker hypothesis:

\begin{theorem}[Wright]\label{thm:wright}
Assume the Heath-Brown conjecture on least primes in arithmetic progressions. Then
\[
C(x) \gg x^{1 - R(x)}, \quad \text{where } R(x) = (2 + o(1)) \frac{\log \log \log \log x}{\log \log \log x}.
\]
\end{theorem}

The Heath-Brown conjecture is implied by Elliott-Halberstam-type hypotheses but is weaker than the full Elliott-Halberstam conjecture.

\subsection{Larsen's Breakthroughs}

Larsen has obtained several remarkable unconditional results on Carmichael numbers in recent years:

\begin{theorem}[Larsen, 2022~\cite{larsen2023}]
\label{thm:larsen-short}
For any constant $C > 1/2$, the interval $[x, x + x/(\log x)^C]$ contains at least one Carmichael number for all sufficiently large $x$.
\end{theorem}

This is analogous to Bertrand's postulate for primes. Note that the interval length $x/(\log x)^C = x^{1 - o(1)}$ tends to $x$ as $x \to \infty$, so this does \emph{not} resolve the question of whether intervals $[x, x + x^\theta]$ for fixed $\theta < 1$ always contain Carmichael numbers.

\begin{theorem}[Larsen, 2025~\cite{larsen2025ap}]
\label{thm:larsen-ap}
Every arithmetic progression either contains infinitely many Carmichael numbers or contains none. There is a simple criterion to determine which case holds.
\end{theorem}

\begin{theorem}[Larsen-Wright, 2025~\cite{larsenwright2025}]
\label{thm:larsen-wright}
For every sufficiently large integer $R$, there exists a Carmichael number with exactly $R$ prime factors.
\end{theorem}

This is progress toward the Granville-Pomerance conjecture that for every $k \ge 3$, there exist infinitely many Carmichael numbers with exactly $k$ prime factors, with counting function $C_k(x) = x^{1/k + o_k(1)}$.

\subsection{Computational Tabulation}

Shallue and Webster~\cite{shalluewebster2024} achieved a computational milestone:

\begin{theorem}[Shallue-Webster, 2024]
\label{thm:shallue-webster}
All Carmichael numbers up to $10^{22}$ have been computed. There are exactly $49{,}679{,}870$ such numbers.
\end{theorem}

This extends earlier tabulations by Pinch (up to $10^{21}$) and provides the dataset for the computational analysis in Part~II of this paper.

\section{Open Problems}\label{sec:open}

\subsection{Erdos's Density Conjecture}

Erdos conjectured:

\begin{conjecture}[Erdos]
\label{conj:erdos}
$C(x) = x^{1 - o(1)}$.
\end{conjecture}

The gap between the lower bound $x^{0.3389}$ (Harman) and the conjectured $x^{1-o(1)}$ is enormous. Granville and Pomerance~\cite{granvillepomerance2001} analyzed heuristics for Carmichael number density, identifying two competing models:
\begin{itemize}
\item \textbf{Erdos model:} Based on smooth numbers, predicts $C(x) \sim x \exp(-c (\log x)(\log \log \log x) / (\log \log x))$ for some $c > 0$.
\item \textbf{Pomerance model:} Based on random factorizations, predicts $C(x) \sim x \exp(-c \sqrt{(\log x)(\log \log x)})$, which is much sparser.
\end{itemize}

Determining which model (if either) is correct remains open.

\subsection{Short Interval Density}

\begin{problem}
For fixed $\theta \in (1/2, 1)$, does every interval $[x, x + x^\theta]$ contain at least one Carmichael number for sufficiently large $x$?
\end{problem}

Larsen's Theorem~\ref{thm:larsen-short} does not answer this, as his intervals have length $x^{1-o(1)}$, not $x^\theta$ for fixed $\theta < 1$.

\subsection{Wieferich Primes}

\begin{definition}
A prime $p$ is a \emph{Wieferich prime} if $p^2 \mid 2^{p-1} - 1$.
\end{definition}

Only two Wieferich primes are known: 1093 and 3511. It is unknown whether infinitely many exist. Silverman showed the $abc$ conjecture implies infinitely many non-Wieferich primes. Recent work~\cite{fellini_murty2025} connects number field analogues to the Ankeny-Artin-Chowla conjecture.

\subsection{The PSW Conjecture}

\begin{conjecture}[Baillie-PSW]
\label{conj:psw}
No composite number is simultaneously a strong pseudoprime to base 2 and a Lucas pseudoprime with the standard parameter choice.
\end{conjecture}

A \$620 prize stands for a counterexample or proof. Extensive computation has found none up to $2^{80}$.

\subsection{Carmichael's Totient Conjecture}

\begin{conjecture}[Carmichael, 1907]
No value of Euler's totient function $\varphi(n)$ is attained exactly once.
\end{conjecture}

Ford showed that if a counterexample exists, a positive proportion of integers are also counterexamples. The conjecture remains open.

\section{Computational Analysis of Carmichael Number Distribution}\label{sec:computational}

We now turn to original computational analysis using the Shallue-Webster dataset of all 49,679,870 Carmichael numbers up to $10^{22}$.

\subsection{Data and Methodology}

The analysis uses the Carmichael number counts $C(10^k)$ for $k = 3, 4, \ldots, 22$ from the OEIS sequence A055553 and the Shallue-Webster tabulation. Table~\ref{tab:counts} presents the data:

\begin{table}[ht]
\centering
\caption{Carmichael number counts and empirical exponents}\label{tab:counts}
\begin{tabular}{rrrl}
\toprule
$x$ & $C(x)$ & $\alpha(x) = \log C(x) / \log x$ & $\log \log x$ \\
\midrule
$10^{6}$ & 43 & 0.272 & 2.77 \\
$10^{9}$ & 646 & 0.312 & 3.10 \\
$10^{12}$ & 8,241 & 0.326 & 3.36 \\
$10^{15}$ & 105,212 & 0.335 & 3.56 \\
$10^{18}$ & 1,401,644 & 0.341 & 3.74 \\
$10^{21}$ & 20,138,200 & 0.347 & 3.89 \\
$10^{22}$ & 49,679,870 & 0.350 & 3.93 \\
\bottomrule
\end{tabular}
\end{table}

We exclude $C(10^3) = 1$ from regression analysis as it is an outlier driven by the single Carmichael number 561 below $10^3$.

\subsection{Effective Exponent Function}

Define the \emph{effective exponent}
\[
\alpha(x) = \frac{\log C(x)}{\log x}.
\]
If Erdos's conjecture holds, $\alpha(x) \to 1$ as $x \to \infty$. The data shows $\alpha(x)$ increasing monotonically from 0.272 at $10^6$ to 0.350 at $10^{22}$, supporting the conjecture. The rate of convergence is very slow.

\subsection{Testing the Erdos-Pomerance Heuristic}

The Erdos model predicts
\[
\alpha(x) = 1 - (d + o(1)) \frac{\log \log \log x}{\log \log x}
\]
for some constant $d > 0$. We fit this model by computing
\[
d(x) = \frac{(1 - \alpha(x)) \log \log x}{\log \log \log x}.
\]

Table~\ref{tab:erdos-fit} presents the fitted values:

\begin{table}[ht]
\centering
\caption{Fitting the Erdos-Pomerance constant}\label{tab:erdos-fit}
\begin{tabular}{rl}
\toprule
$x$ & $d(x)$ \\
\midrule
$10^{6}$ & 1.98 \\
$10^{9}$ & 1.88 \\
$10^{12}$ & 1.86 \\
$10^{15}$ & 1.86 \\
$10^{18}$ & 1.87 \\
$10^{21}$ & 1.87 \\
$10^{22}$ & 1.87 \\
\bottomrule
\end{tabular}
\end{table}

The constant $d$ stabilizes near $1.87$ for $x \ge 10^{12}$, providing strong empirical support for the Erdos-Pomerance form with $d \approx 1.87$.

\subsection{Comparison with Alternative Models}

We also tested the functional form $\alpha(x) = 1 - A / (\log \log x)^B$, fitting $A$ and $B$ by nonlinear regression. However, the constant $A$ drifts monotonically from 2.2 to 3.6 across the data range, indicating this model does not fit the data as well as the Erdos-Pomerance form. The narrow range of $\log \log x$ (from 2.77 to 3.93, a 42\% variation) limits the discriminatory power of regression over this interval.

\subsection{Distribution of Prime Factor Counts}

For each Carmichael number $n$ in the dataset, let $\omega(n)$ denote the number of distinct prime factors. The AGP construction produces Carmichael numbers with many prime factors (growing with $\log x$), while the smallest Carmichael numbers tend to have $\omega(n) = 3$ (the minimum possible by Korselt's criterion).

Analysis of the Shallue-Webster data (to the extent the full factorizations are accessible) would reveal how $\omega(n)$ is distributed as a function of $n$. The Granville-Pomerance heuristic predicts ``typical'' Carmichael numbers near $x$ have $\omega(n) \sim (\log x)^{1-\epsilon}$ for small $\epsilon > 0$, but the precise distribution remains to be measured.

\subsection{Smoothness Profile}

For a Carmichael number $n = p_1 \cdots p_r$, Korselt's criterion requires $(p_i - 1) \mid (n-1)$ for each $i$. Define the \emph{smoothness parameter}
\[
s(n) = \max_{i=1,\ldots,r} \frac{\log P^+(p_i - 1)}{\log n},
\]
where $P^+(m)$ denotes the largest prime factor of $m$. This measures how ``smooth'' the values $p_i - 1$ are relative to $n$.

In the AGP construction with smoothness parameter $y = x^{1/7}$, the primes $p$ are chosen with $P^+(p-1) \le y$, giving $s(n) \sim 1/7 \approx 0.143$ for AGP-constructed Carmichael numbers. If ``wild'' Carmichael numbers (not arising from AGP-type constructions) exist in abundance with different smoothness profiles, this would be reflected in the distribution of $s(n)$.

Computing this distribution requires factoring $p_i - 1$ for each prime factor $p_i$ of each Carmichael number in the dataset---a computationally intensive but feasible task for $n \le 10^{22}$ (since $p_i \le 10^{22}$ and standard factorization methods suffice).

\section{Conjectures}\label{sec:conjectures}

Based on the computational analysis, we propose two conjectures with explicit, testable predictions.

\subsection{Refined Erdos Conjecture}

\begin{conjecture}\label{conj:refined-erdos}
The Carmichael number counting function satisfies
\[
C(x) = x^{1 - (1.87 + o(1)) \frac{\log \log \log x}{\log \log x}}
\]
where the error term $o(1)$ is bounded by $|o(1)| \le B / (\log \log x)^{1/2}$ for some explicit constant $B > 0$ and all sufficiently large $x$.
\end{conjecture}

This calibrates the Erdos-Pomerance heuristic to the Shallue-Webster data. The constant $1.87$ is determined empirically from Table~\ref{tab:erdos-fit}. The conjecture predicts:
\begin{align*}
C(10^{25}) &\approx 10^{25} \cdot 10^{-7.75} \approx 1.8 \times 10^{17}, \\
C(10^{30}) &\approx 10^{30} \cdot 10^{-9.08} \approx 8.3 \times 10^{20}.
\end{align*}

These predictions can be tested against future tabulations.

\subsection{Short Interval Density}

The following conjecture extends the Erdos density to short intervals:

\begin{conjecture}\label{conj:short-interval}
For fixed $\theta \in (1/2, 1)$, the density of Carmichael numbers in intervals $[x, x + x^\theta]$ satisfies
\[
C(x, x + x^\theta) \ge x^{\theta - (1.87 + o(1)) \frac{\log \log \log x}{\log \log x}}
\]
for all sufficiently large $x$.
\end{conjecture}

Under a ``uniform distribution'' heuristic, Carmichael numbers near $x$ have density approximately $C(x)/x \approx x^{\alpha(x) - 1}$, so an interval of length $x^\theta$ should contain roughly
\[
C(x, x + x^\theta) \sim x^\theta \cdot x^{\alpha(x) - 1} = x^{\theta + \alpha(x) - 1}.
\]
Conjecture~\ref{conj:short-interval} asserts this heuristic holds (at least as a lower bound) with the Erdos-Pomerance form for $\alpha(x)$.

For $\theta = 0.9$ and $x = 10^{25}$, the conjecture predicts
\[
C(10^{25}, 10^{25} + 10^{22.5}) \ge 10^{22.5} \cdot 10^{-7.75} \approx 5.6 \times 10^{14}.
\]

\subsection{Testability and Limitations}

Both conjectures are testable against future computational tabulations. However, we note several limitations:
\begin{itemize}
\item The range $x \le 10^{22}$ may be insufficient to distinguish the Erdos model from the Pomerance model, as both predict very similar counts at this scale.
\item The constant $1.87$ in Conjecture~\ref{conj:refined-erdos} is fitted from only 7 data points with $\log \log x$ varying over a narrow range (2.77 to 3.93). Confidence intervals are large.
\item Conjecture~\ref{conj:short-interval} assumes a ``uniform distribution'' heuristic that may not hold if Carmichael numbers exhibit clustering or other local irregularities.
\end{itemize}

Despite these caveats, the conjectures provide concrete, falsifiable predictions that extend beyond the current computational frontier.

\section{Discussion}

Fermat's Little Theorem, now 386 years old, continues to generate profound consequences and open problems. The recent breakthroughs by Larsen and Larsen-Wright on Carmichael numbers in short intervals and arithmetic progressions, combined with the Shallue-Webster computational tabulation, have dramatically advanced our understanding of pseudoprime distribution. Yet the gap between the unconditional lower bound $C(x) \ge x^{0.3389}$ and the conjectured $C(x) = x^{1-o(1)}$ remains vast.

The computational analysis in this paper supports the Erdos-Pomerance heuristic over alternative models, with an empirically determined constant near $d \approx 1.87$. Whether this constant is universal or whether it drifts at larger $x$ remains to be seen. The proposed conjectures (Conjectures~\ref{conj:refined-erdos} and~\ref{conj:short-interval}) provide testable predictions for the next generation of tabulations.

Other major open problems---the infinitude of Wieferich primes, the PSW conjecture, Carmichael's totient conjecture---remain as tantalizing as ever. And the advent of quantum computing threatens to upend the cryptographic edifice built on FLT, driving research into post-quantum alternatives.

Yet the elegance and power of Fermat's original observation endures: $a^{p-1} \equiv 1 \pmod{p}$. Few theorems in mathematics can claim such a span from classical number theory to modern cryptography, from pure combinatorics to computational complexity. Fermat's Little Theorem is indeed a giant upon whose shoulders much of modern mathematics stands.

\bibliographystyle{amsplain}
\bibliography{fermat_carmichael}
2nd ed., Springer, 2005.

\bibitem{diffiehellman1976}
Whitfield Diffie and Martin E. Hellman,
\emph{New directions in cryptography},
IEEE Transactions on Information Theory \textbf{22} (1976), no.~6, 644--654.

\bibitem{euler1736}
Leonhard Euler,
\emph{Theorematum quorundam ad numeros primos spectantium demonstratio},
Commentarii Academiae Scientiarum Petropolitanae \textbf{8} (1736), 141--146.

\bibitem{euler1763}
Leonhard Euler,
\emph{Theoremata circa residua ex divisione potestatum relicta},
Novi Commentarii Academiae Scientiarum Petropolitanae \textbf{7} (1763), 49--82.

\bibitem{fellini_murty2025}
Fellini and M. Ram Murty,
\emph{Wieferich primes in number fields and the conjectures of {A}nkeny--{A}rtin--{C}howla and {M}ordell},
arXiv preprint arXiv:2508.08472, 2025.

\bibitem{granvillepomerance2001}
Andrew Granville and Carl Pomerance,
\emph{Two contradictory conjectures concerning {C}armichael numbers},
Mathematics of Computation \textbf{71} (2002), no.~238, 883--908.

\bibitem{harman2005}
Glyn Harman,
\emph{On the number of Carmichael numbers up to $x$},
Bulletin of the London Mathematical Society \textbf{37} (2005), 641--650.

\bibitem{harman2008}
Glyn Harman,
\emph{Watt's mean value theorem and Carmichael numbers},
International Journal of Number Theory \textbf{4} (2008), no.~2, 241--248.

\bibitem{hernandez2020}
de Melo Hern\'{a}ndez, Hern\'{a}ndez Melo, and Tapia-Recillas,
\emph{Fermat's {L}ittle {T}heorem and {E}uler's {T}heorem in a class of rings},
arXiv preprint arXiv:2012.06949, 2020.

\bibitem{isaacs_pournaki2005}
I. Martin Isaacs and M. R. Pournaki,
\emph{Generalizations of {F}ermat's {L}ittle {T}heorem via group theory},
American Mathematical Monthly \textbf{112} (2005), 734--740.

\bibitem{korselt1899}
Alwin Korselt,
\emph{Probl\`{e}me chinois},
L'interm\'{e}diaire des math\'{e}maticiens \textbf{6} (1899), 142--143.

\bibitem{larsen2023}
Daniel Larsen,
\emph{Bertrand's postulate for {C}armichael numbers},
International Mathematics Research Notices \textbf{2023} (2023), no.~15, 13072--13098.

\bibitem{larsen2025ap}
Daniel Larsen,
\emph{Carmichael numbers in all possible arithmetic progressions},
arXiv preprint arXiv:2504.09056, 2025.

\bibitem{larsenwright2025}
Daniel Larsen and Thomas Wright,
\emph{Carmichael numbers with a specified number of prime factors},
arXiv preprint arXiv:2510.16632, 2025.

\bibitem{rabin1980}
Michael O. Rabin,
\emph{Probabilistic algorithm for testing primality},
Journal of Number Theory \textbf{12} (1980), no.~1, 128--138.

\bibitem{rsa1978}
Ronald L. Rivest, Adi Shamir, and Leonard Adleman,
\emph{A method for obtaining digital signatures and public-key cryptosystems},
Communications of the ACM \textbf{21} (1978), no.~2, 120--126.

\bibitem{shalluewebster2024}
Andrew Shallue and Jonathan Webster,
\emph{Advances in tabulating {C}armichael numbers},
Research in Number Theory \textbf{10} (2024), article 8.

\bibitem{shor1994}
Peter W. Shor,
\emph{Algorithms for quantum computation: discrete logarithms and factoring},
Proceedings of the 35th Annual Symposium on Foundations of Computer Science, IEEE, 1994, pp.~124--134.

\bibitem{washington2008}
Lawrence C. Washington,
\emph{Elliptic Curves: Number Theory and Cryptography},
2nd ed., Chapman and Hall/CRC, 2008.

\bibitem{wright2020}
Thomas Wright,
\emph{A conditional density for {C}armichael numbers},
Bulletin of the Australian Mathematical Society \textbf{101} (2020), 379--388.

\end{thebibliography}

\end{document}
